\subsection{SGD and DoG}\label{sec:method-m1}
\owner{Member 1}

Plain SGD uses $\theta_{t+1}=\theta_t-\eta_tg_t$. Its learning rate must balance progress
and stability. DoG instead constructs a scalar step size from distance traveled from the
initial point and accumulated gradient norms \citep{ivgi2023dog}.
We will test whether this reduces sensitivity to an externally selected learning rate.
The paper's stochastic convex guarantee concerns a specified variant and its assumptions;
it is not a guarantee for arbitrary nonconvex training. Conditioning diagnostics probe the
limits of the mechanism rather than assuming that scalar adaptation removes anisotropy.

\fillin{Member 1: Write the exact DoG update and initialization, distinguish it from the theoretical variant, and specify the chosen distance floor. Compare with tuned SGD. Report the parameter-state count, tuning evaluations, and correctness checks.}
